% Person 5's report section: "Structured memory: design, final memory, and a
% trajectory where memory helped" (~1 page). Self-contained fragment; \input
% this into the main report.
%
% PACKAGE NEEDED: listings (for \begin{lstlisting} with breaklines, used to
% quote verbatim memory/snapshot text without it overflowing the page width).
% Add \usepackage{listings} to the preamble. No other package is required.

\section{Structured memory: design, final memory, and a trajectory where memory helped}

\paragraph{Design.} \texttt{StructuredMemory} (\texttt{task5\_structured\_memory/structured.py})
is deterministic and makes zero LLM calls, so it adds no API cost beyond the agent's own
calls. It starts empty; nothing about C1/C2/E1/E2 is pre-written anywhere in the code.

\emph{Write.} A rule is created only when the checker's reviewer feedback on a failed
attempt is classified, by substring match, into one of the four conventions. Its
\texttt{text} field is stored \emph{verbatim} as that reviewer message -- never paraphrased.

\emph{Retrieve.} The day rule (C2) is withheld for whole-month tasks and shown only for
single-day-scoped questions -- a retrieval heuristic, not a fact the agent was ever told.
Station-based exceptions (Wingra) are tagged \emph{applies to this task} or \emph{does not
apply to this task} against the current task's station; the per-variable exception (E2) is
shown neutrally, since its reviewer wording never names a station or variable to check.

\emph{Update.} When feedback names an exception, it edits the parent rule in place (attaching
the exception and narrowing its scope) rather than appending a second, possibly-contradictory
rule. A single reviewer message can name more than one convention at once (Person 1's
``both\_rules'' message); \texttt{classify\_feedback} returns every kind such a message
names, and \texttt{update} writes or edits every rule it names, storing the same verbatim text
as each rule's \texttt{text}. Each rule's scope is assembled from evidence -- the
station/variable/day-type of the task that first taught it, plus later tasks where it was
active and the agent solved the task on the first attempt, labeled \emph{retrieved, solved
first try}. That label is deliberately weak: a first-attempt success does not prove the agent
used the memory, only that it didn't need correcting.

\emph{Forget.} The first time an exception is attached to a rule, the pre-exception version
of that rule is frozen into a separate, timestamped \emph{retired} entry -- visible in
\texttt{snapshot()} but never returned by \texttt{retrieve()}. Separately, if a rule was handed
to the agent and the task still failed on feedback naming that same rule, it is logged under
\texttt{retrieved\_but\_not\_applied} on the rule (count and task IDs) -- direct evidence of
memory being retrieved and not applied (see the discussion notes for a worked example).

Storing scope \emph{with} each rule, rather than the rule alone, is what T5, T6, T8, and T9 are
designed to test: a memory that only remembers ``shift the day to 08:00'' or ``drop flagged
readings'' has no way to except Wingra or per-variable filtering without overwriting itself
into a contradiction. Recording \emph{where} a rule was learned and where its exceptions were
discovered lets the exception edit the rule precisely instead of replacing it.

\paragraph{Final structured memory.} All claims and numbers below come from the full protocol
run (three repeats, all three conditions) that landed in \texttt{outputs/} under run
\texttt{20261013-200839}. Two earlier single-pass structured-only runs
(\texttt{20261013-182822}, \texttt{20261013-193556}) also exist, but they ran against an
earlier version of the task data (e.g. T7's ground-truth answer has since changed) and are not
counted in any number in this section. Below is the structured-memory snapshot after T10 of
\emph{repeat~1} of \texttt{20261013-200839}, chosen because in this repeat the Wingra exception
(E1) was learned at T5, the task the sequence was designed to introduce it at (see below for
why the other repeats differ). The \texttt{[C1]} and \texttt{[C2]} active rules are reproduced
in full; the one retired entry is included as-is (it is short enough not to need summarizing).

{\footnotesize
\begin{lstlisting}[breaklines=true, columns=fullflexible]
Structured memory snapshot
===========================

[C1] status=active
  Rule (verbatim feedback): Readings aren't trustworthy unless their own quality flag is 0; drop estimated, suspect, and offline values.
  Scope: learned at monona/water_temp/single_day (T2); retrieved, solved first try at monona/dissolved_oxygen/single_day (T4); retrieved, solved first try at mendota/wind_speed/whole_month (T6); retrieved, solved first try at mendota+monona/water_temp/single_day (T7); retrieved, solved first try at mendota+wingra/water_temp/single_day (T8); retrieved, solved first try at wingra/wind_speed/single_day (T9); retrieved, solved first try at monona/dissolved_oxygen/single_day (T10)

[C2] status=active
  Rule (verbatim feedback): Lab days run 08:00 to 08:00 local for the v1 stations, not midnight to midnight.
  Scope: learned at monona/water_temp/single_day (T2); retrieved, solved first try at monona/dissolved_oxygen/single_day (T4); retrieved, solved first try at mendota+monona/water_temp/single_day (T7); retrieved, solved first try at monona/dissolved_oxygen/single_day (T10)
  Exception (station=wingra): Wingra's v2 logger already records on the standard calendar day; don't apply the 08:00 observation-day shift. [seen at: T5]
  Retrieved but not applied: 1 time(s) (T3)

[C2@T5] status=retired superseded_by=C2 retired_at=T5
  Reason: narrowed by exception learned at T5: Wingra's v2 logger already records on the standard calendar day; don't apply the 08:00 observation-day shift.
  Rule (verbatim feedback): Lab days run 08:00 to 08:00 local for the v1 stations, not midnight to midnight.
  Scope: learned at monona/water_temp/single_day (T2); retrieved, solved first try at monona/dissolved_oxygen/single_day (T4)
  Retrieved but not applied: 1 time(s) (T3)
\end{lstlisting}
}

Two things are notably \emph{inconsistent} across repeats of this same run, and we report that
plainly rather than smooth it over. First, \emph{when} E1 is learned varies: in repeat~1
(shown above) it was learned at T5 as intended, but in repeats~0 and~2 the agent solved T5 on
the first attempt and the Wingra exception was instead learned later, at T8 (both repeats'
first T8 attempt got answer $-0.48$ and the identical ``Wingra's v2 logger...'' feedback before
correcting to $0.08$). Across the three repeats of this run, E1 fired at T5 once (repeat~1) and
at T8 twice (repeats~0 and~2) -- never consistently where the task sequence intends. Second,
C1's per-variable exception (E2, planted for T6/T9) has \emph{never} fired in any of the three
repeats -- T6 was solved on the first attempt every time, so the row-drop mistake it is meant
to catch has simply never happened, and \texttt{structured.py}'s exception-editing logic for C1
remains untested by real data.

\paragraph{Trajectory where memory helped.} T7 (``Mendota minus Monona, July 20'', needing both
C1 and C2) is the most robust example across run \texttt{20261013-200839}: structured solved it
on the \emph{first} attempt in all three repeats (1 attempt, 4 model calls each), while the
\texttt{none} baseline needed more than one attempt in all three repeats (3, 2, and 2 attempts).
In repeat~0, structured received this memory text and solved it immediately:

{\footnotesize
\begin{lstlisting}[breaklines=true, columns=fullflexible]
1. Readings aren't trustworthy unless their own quality flag is 0; drop estimated, suspect, and offline values.
   Scope: learned at mendota/water_temp/whole_month (T1); retrieved, solved first try at monona/dissolved_oxygen/single_day (T4); retrieved, solved first try at wingra/water_temp/single_day (T5); retrieved, solved first try at mendota/wind_speed/whole_month (T6)
2. Lab days run 08:00 to 08:00 local for the v1 stations, not midnight to midnight.
   Scope: learned at monona/water_temp/single_day (T2); retrieved, solved first try at monona/dissolved_oxygen/single_day (T4); retrieved, solved first try at wingra/water_temp/single_day (T5)
\end{lstlisting}
}

With that memory, the agent filtered both stations to \texttt{water\_temp\_qc==0} and windowed
each day 08:00--08:00, answering $0.57$ on its single attempt. The \texttt{none} baseline, with
no memory, took three attempts on repeat~0 of the same run: its first answer ($0.11$) was
generic-wrong, its second ($-0.43$) drew the exact reviewer message ``Lab days run 08:00 to
08:00 local for the v1 stations, not midnight to midnight,'' and only the third attempt matched.
Repeats 1 and 2 of \texttt{none} each needed two attempts, always including that same C2
reviewer message before correcting -- the cold agent had to rediscover the day rule on T7 itself
every time. T10 shows the same pattern even more starkly: structured was first-try in all three
repeats, while \texttt{none} never answered T10 correctly within three attempts in two of the
three repeats (final \texttt{success=False} for both).

T8 -- suggested as a candidate trajectory -- turned out \emph{not} to be a clean ``memory
helped'' example once checked against the data, so it is not used as the main trajectory here.
Across the three repeats of \texttt{20261013-200839}, T8 was the task where E1 was first
learned (via a failed first attempt) in two of them (repeats~0 and~2); only in repeat~1, where
E1 was already known from T5, did T8 succeed on the first attempt with the rule already in
memory. T8 is therefore mostly a ``memory learned an exception'' example, not a ``memory
helped'' one, in this data.
